% GENERATED FILE. Edit canonical lessons, then run npm run generate:books.
% canonical-source-hash: fnv1a64:613459451e140fb1
% canonical-lessons: ZH-C296-yiwai, ZH-C296-chehuo, ZH-C296-shache, ZH-C296-xinshou, ZH-C296-peiqian

\chapter{A Crash}
\label{ch:zh-a-crash}
\hlchaptermodality{296}

\begin{chapteropening}
\textbf{By the end of this chapter:} \emph{I can talk about a crash on the road: the brakes, a new driver and paying for the damage.}

Complete the last lesson of chapter 296: I can talk about a crash on the road: the brakes, a new driver and paying for the damage.
\end{chapteropening}

\section[\texorpdfstring{\zh{意外} (yìwài)}{意外 (yìwài)}]{\zh{意外} (yìwài) — mishap; something unexpected}

\label{lesson:ZH-C296-yiwai}



\begin{quote}
Before the new one: say the Mandarin for after the event, then the Mandarin for right there on the spot.
\end{quote}

\subsection*{You'll want to know: \zh{意外}}
\textbf{\zh{意外}} — \emph{yìwài} — \textquotedblleft{}mishap; something unexpected\textquotedblright{}.

\begin{quote}
\textbf{\zh{意外发生的时候我在家}} \emph{yìwài fāshēng de shíhou, wǒ zài jiā} — When the accident happened, I was at home.
\end{quote}

\begin{cousinweb}
\textbf{\zh{意外}} (\emph{yìwài}) is built from \zh{意} 'an intention' and \zh{外} 'outside', both inherited from Middle Chinese, \emph{'iH} and \emph{ngwajH} (Guangyun, in Baxter's transcription), which glosses the first 'will' and the second 'outside' and 'far': outside what anyone intended.
\end{cousinweb}

\begin{grammarlens}[title={What you've built}]
One more word for telling what happened.
\end{grammarlens}

\subsection*{Your turn}
\begin{itemize}
\raggedright
  \item \emph{Say it:} \emph{yìwài}
  \item \emph{Say it:} \emph{yìwài}, once more
  \item \emph{Say it:} \emph{yìwài fāshēng de shíhou, wǒ zài jiā} — When the accident happened, I was at home.
  \item \emph{Recall:} say the Mandarin for after the event, then the Mandarin for right there on the spot, then say \emph{yìwài} again
\end{itemize}

\begin{tcolorbox}[breakable,colback=teal!4,colframe=teal!35!black,title={Before you move on}]
What does \zh{意外} mean? (\textquotedblleft{}Mishap; something unexpected\textquotedblright{}.) Say it once more.
\end{tcolorbox}
\section[chēhuò]{chēhuò — traffic accident; car crash}

\label{lesson:ZH-C296-chehuo}



\begin{quote}
Before the new one: say the Mandarin for right there on the spot, then the Mandarin for something unexpected.
\end{quote}

\subsection*{You'll want to know: chēhuò}
\emph{chēhuò} — \textquotedblleft{}traffic accident; car crash\textquotedblright{}.

\begin{quote}
\emph{tā zài chēhuò lǐ shòushāng le} — He was hurt in a car crash.
\end{quote}

Its characters come later; for now it is a word for the ear.

\begin{cousinweb}
\emph{chēhuò} is built from \zh{车} 'a vehicle' and \emph{huò} 'a disaster', both inherited from Middle Chinese, \emph{tsyhae} and \emph{hwaX} (Guangyun, in Baxter's transcription), which glosses the second 'harm'.
\end{cousinweb}

\begin{grammarlens}[title={What you've built}]
One more word for telling what happened.
\end{grammarlens}

\subsection*{Your turn}
\begin{itemize}
\raggedright
  \item \emph{Say it:} \emph{chēhuò}
  \item \emph{Say it:} \emph{chēhuò}, once more
  \item \emph{Say it:} \emph{tā zài chēhuò lǐ shòushāng le} — He was hurt in a car crash.
  \item \emph{Recall:} say the Mandarin for right there on the spot, then the Mandarin for something unexpected, then say \emph{chēhuò} again
\end{itemize}

\begin{tcolorbox}[breakable,colback=teal!4,colframe=teal!35!black,title={Before you move on}]
What does chēhuò mean? (\textquotedblleft{}Traffic accident; car crash\textquotedblright{}.) Say it once more.
\end{tcolorbox}
\section[shāchē]{shāchē — to put on the brakes; a brake}

\label{lesson:ZH-C296-shache}



\begin{quote}
Before the new one: say the Mandarin for something unexpected, then the Mandarin for a car crash.
\end{quote}

\subsection*{You'll want to know: shāchē}
\emph{shāchē} — \textquotedblleft{}to put on the brakes; a brake\textquotedblright{}.

\begin{quote}
\emph{kuài shāchē, bùrán jiù zhuàng le} — Brake, quick, or we'll hit it!
\end{quote}

Its characters come later; for now it is a word for the ear.

\begin{cousinweb}
\emph{shāchē} is built from \emph{shā} 'to stop' and \zh{车} 'a vehicle'. The second is inherited from Middle Chinese \emph{tsyhae} (Guangyun, in Baxter's transcription). The Guangyun lists the first as \emph{tsrhaet}, 'the pillar of a temple', the source of its other reading chà; the shā of 'to brake' has no source we can open, so its history is not settled.
\end{cousinweb}

\begin{grammarlens}[title={What you've built}]
One more word for telling what happened.
\end{grammarlens}

\subsection*{Your turn}
\begin{itemize}
\raggedright
  \item \emph{Say it:} \emph{shāchē}
  \item \emph{Say it:} \emph{shāchē}, once more
  \item \emph{Say it:} \emph{kuài shāchē, bùrán jiù zhuàng le} — Brake, quick, or we'll hit it!
  \item \emph{Recall:} say the Mandarin for something unexpected, then the Mandarin for a car crash, then say \emph{shāchē} again
\end{itemize}

\begin{tcolorbox}[breakable,colback=teal!4,colframe=teal!35!black,title={Before you move on}]
What does shāchē mean? (\textquotedblleft{}To put on the brakes; a brake\textquotedblright{}.) Say it once more.
\end{tcolorbox}
\section[\texorpdfstring{\zh{新手} (xīnshǒu)}{新手 (xīnshǒu)}]{\zh{新手} (xīnshǒu) — novice; beginner}

\label{lesson:ZH-C296-xinshou}



\begin{quote}
Before the new one: say the Mandarin for a car crash, then the Mandarin for to put on the brakes.
\end{quote}

\subsection*{You'll want to know: \zh{新手}}
\textbf{\zh{新手}} — \emph{xīnshǒu} — \textquotedblleft{}novice; beginner\textquotedblright{}.

\begin{quote}
\emph{wǒ shì xīnshǒu, kāichē hěn màn} — I'm a beginner, so I drive slowly.
\end{quote}

\begin{cousinweb}
\textbf{\zh{新手}} (\emph{xīnshǒu}) is built from \zh{新} 'new' and \zh{手} 'a hand', both inherited from Middle Chinese, \emph{sin} and \emph{syuwX} (Guangyun, in Baxter's transcription): a new hand, as in English.
\end{cousinweb}

\begin{grammarlens}[title={What you've built}]
One more word for telling what happened.
\end{grammarlens}

\subsection*{Your turn}
\begin{itemize}
\raggedright
  \item \emph{Say it:} \emph{xīnshǒu}
  \item \emph{Say it:} \emph{xīnshǒu}, once more
  \item \emph{Say it:} \emph{wǒ shì xīnshǒu, kāichē hěn màn} — I'm a beginner, so I drive slowly.
  \item \emph{Recall:} say the Mandarin for a car crash, then the Mandarin for to put on the brakes, then say \emph{xīnshǒu} again
\end{itemize}

\begin{tcolorbox}[breakable,colback=teal!4,colframe=teal!35!black,title={Before you move on}]
What does \zh{新手} mean? (\textquotedblleft{}Novice; beginner\textquotedblright{}.) Say it once more.
\end{tcolorbox}
\section[péiqián]{péiqián — to pay for damage}

\label{lesson:ZH-C296-peiqian}



\begin{quote}
Before the new one: say the Mandarin for to put on the brakes, then the Mandarin for a beginner at the wheel.
\end{quote}

\subsection*{You'll want to know: péiqián}
\emph{péiqián} — \textquotedblleft{}to pay for damage\textquotedblright{}.

\begin{quote}
\emph{tā zhuàng le wǒ de chē, děi péiqián} — He hit my car, so he has to pay for it.
\end{quote}

Its characters come later; for now it is a word for the ear.

\begin{cousinweb}
\emph{péiqián} is built from \emph{péi} 'to make good a loss' and \zh{钱} 'money'. The second is inherited from Middle Chinese \emph{dzjen} (Guangyun, in Baxter's transcription); the Guangyun has no entry for the first, so its history is not settled.
\end{cousinweb}

\begin{grammarlens}[title={What you've built}]
One more word for telling what happened.
\end{grammarlens}

\subsection*{Your turn}
\begin{itemize}
\raggedright
  \item \emph{Say it:} \emph{péiqián}
  \item \emph{Say it:} \emph{péiqián}, once more
  \item \emph{Say it:} \emph{tā zhuàng le wǒ de chē, děi péiqián} — He hit my car, so he has to pay for it.
  \item \emph{Recall:} say the Mandarin for to put on the brakes, then the Mandarin for a beginner at the wheel, then say \emph{péiqián} again
\end{itemize}

\begin{tcolorbox}[breakable,colback=teal!4,colframe=teal!35!black,title={Before you move on}]
What does péiqián mean? (\textquotedblleft{}To pay for damage\textquotedblright{}.) Say it once more.
\end{tcolorbox}
